% Part: first-order-logic 
% Chapter: axiomatic-deduction 
% Section: deduction-theorem-quantifiers

\documentclass[../../../include/open-logic-section]{subfiles}

\begin{document}

\olfileid{fol}{axd}{ddq}

\olsection{مسوّراں سمیت استنباط دا قضیہ}

\begin{thm}[استنباط دا قضیہ]
\ollabel{thm:deduction-thm-q} جے $\Gamma \cup \{!A\} \Proves !B$،
تاں $\Gamma \Proves !A \lif !B$۔
\end{thm}

\begin{proof}
اسیں فیر $!B$ دے $\Gamma \cup \{!A\}$ توں !!{derivation} دی
لمبائی اُتے استقرا کردے آں۔

استقرا دی بنیاد دا ثبوت بالکل اوہی اے جیہڑا
\olref[ded]{thm:deduction-thm} دے ثبوت وچ اے۔

استقرائی قدم لئی فیر فرض کرو کہ $!B$ دا !!{derivation}،
$\Gamma \cup \{!A\}$ توں، ایسے قدم~$!B$ اُتے ختم ہوندا اے جیہڑا
استنباط دے قاعدے نال درست ٹھہردا اے۔ جے ایہ قاعدہ اثباتِ مقدم اے،
تاں اسیں \olref[ded]{thm:deduction-thm} دے ثبوت وانگ اگے ودھدے آں۔
جے قاعدہ \QR اے، تاں اسیں جاندے آں کہ $!B \ident !C \lif \lforall[x][!D(x)]$
تے $!C \lif !D(a)$ دے روپ دا !!a{formula} پہلے !!{derivation}
وچ آ چکیا اے؛ ایتھے $a$ نہ~$!C$ وچ، نہ~$!A$ وچ، تے نہ~$\Gamma$
وچ آؤندا اے۔ ایس لئی ساڈے کول اے:
% PNBAXD-005: دوجی میٹا شرطیہ وچ اک کھلا بریکٹ بند کیتا۔
\begin{align*}
  \Gamma \cup \{!A\} & \Proves !C \lif !D(a),\\
  \intertext{تے استقرائی مفروضہ لاگو ہوندا اے، یعنی ساڈے کول اے}
    \Gamma & \Proves !A \lif (!C \lif !D(a)).\\
  \intertext{ایس توں}
  & \Proves (!A \lif (!C \lif !D(a))) \lif ((!A \land !C) \lif !D(a))\\
  \intertext{تے قاعدۂ اثباتِ مقدم نال ملدا اے}
  \Gamma & \Proves (!A \land !C) \lif !D(a).\\
  \intertext{کیوں جے خاص مستقل دی غیر حاضری والی شرط ہُن وی پوری اے، اسیں ایس !!{derivation} وچ \QR نال درست ہون والا اک قدم ودھا سکدے آں، تے حاصل کردے آں}
    \Gamma & \Proves (!A \land !C) \lif \lforall[x][!D(x)].\\
    \intertext{ساڈے کول ایہ وی اے}
    & \Proves ((!A \land !C) \lif \lforall[x][!D(x)]) \lif (!A \lif (!C \lif \lforall[x][!D(x)])),\\
    \intertext{لہٰذا قاعدۂ اثباتِ مقدم نال}
    \Gamma & \Proves !A \lif (!C \lif \lforall[x][!D(x)]),
\end{align*}
% PNBAXD-004: آخری نتیجے وچ ماخذ توں رہ گیا A والا شرطیہ بحال کیتا۔
یعنی $\Gamma \Proves !A \lif !B$۔

اوہ حالت جتھے $!B$ قاعدے \QR نال درست ٹھہردا اے، پر اوہدا روپ
$\lexists[x][!D(x)] \lif !C$ اے، مشق لئی چھڈدے آں۔
\end{proof}

\begin{prob}
  \olref[fol][axd][ddq]{thm:deduction-thm-q} دا ثبوت پورا کرو۔
\end{prob}

\end{document}
